% !TeX root = ../tfg.tex
% !TeX encoding = utf8
%
%*******************************************************
% Capítulo 6. Inferencia estadística
%*******************************************************

\chapter{Inferencia estadística}\label{cap:inferencia}

\section{El test de Mann-Whitney}\label{sec:mann-whitney}

% Referencias: Mann y Whitney (1947), Wilcoxon (1945); implementación de SciPy (Virtanen et al., 2020).

\begin{definicion}[Estadístico U de Mann-Whitney]\label{def:estadistico-u}
% Muestras X_1..X_{n_1}, Y_1..Y_{n_2}; U = \sum_i \sum_j \Ind(X_i > Y_j) (+ 1/2 por empate).
\pendiente{definición}
\end{definicion}

\begin{proposicion}[Media y varianza de U bajo la hipótesis nula]\label{prop:u-momentos}
% E[U] = n_1 n_2 / 2; Var(U) = n_1 n_2 (n_1 + n_2 + 1) / 12 sin empates.
\pendiente{enunciado}
\end{proposicion}
\begin{demostracion}
\pendiente{demostración}
\end{demostracion}

\begin{teorema}[Aproximación normal]\label{teo:u-normal}
% z = (U - E[U]) / \sqrt{Var(U)} converge en distribución a N(0, 1).
\pendiente{enunciado}
\end{teorema}

% - Corrección por empates de la varianza.
% - Corrección de continuidad.
% - Saturación: si todas las distancias de un grupo superan a las del otro, U = n_1 n_2 y z
%   alcanza su máximo (38,69 con K = 500 en STOOD-X; empates entre imágenes saturadas).
% - Dependencia: el test supone muestras independientes; en STOOD-X las K^2 distancias de
%   referencia no lo son (vecinos compartidos), lo que afecta a la calibración (capítulo 11).
\begin{observacion}[Empates, continuidad y saturación]\label{obs:u-empates}
\pendiente{desarrollo}
\end{observacion}

\begin{observacion}[Hipótesis de independencia]\label{obs:u-dependencia}
\pendiente{desarrollo}
\end{observacion}

\section{Bootstrap y test de McNemar}\label{sec:bootstrap-mcnemar}

% - Bootstrap no paramétrico: IC por percentiles; bootstrap estratificado (ID y OOD por separado)
%   y bootstrap pareado (mismos índices para dos métodos).
% - Bootstrap por lámina (remuestreo de grupos) frente a por imagen.
% - Test de McNemar para comparar dos clasificadores sobre las mismas imágenes: versión exacta
%   (binomial) y chi-cuadrado con corrección de continuidad.
\begin{definicion}[Intervalo bootstrap por percentiles]\label{def:bootstrap}
\pendiente{definición}
\end{definicion}

\begin{proposicion}[Test de McNemar]\label{prop:mcnemar}
\pendiente{enunciado}
\end{proposicion}

\section{Predicción conforme}\label{sec:prediccion-conforme}

\pendiente{introducción a la predicción conforme}

\endinput
